%% kapa: exam=1 round=37 year=2026 period=1 form=A subject=economics subject_ko=경제학원론 seq=8 number=48
%% source: Q-Net 감정평가사 기출문제 (cst00309) artlSeq=5258924 file=raw/1008/5258924/2026년 제37회 감정평가사 1차 1교시.pdf
%% conversion: kordoc 4.24.51 -> tools/gen1.py (auto); review=proofread-sonnet; flags=none
\begin{question}{48}
모든 면에서 동일한 150개의 개별기업이 존재하는 독점적 경쟁시장에서, 역수요함수와 비용함수는 아래와 같다. 장기균형에서 이 시장에 생존하게 될 최종 기업 수($N$)는? (단, $P$=제품가격, $Q$=개별 기업의 생산량)
\begin{framedtext}
○ 개별 기업의 역수요함수: $P=400-5Q-2N$ \par
○ 개별 기업의 비용함수: $TC=3Q^{2}+40Q+800$ \par
\end{framedtext}
\begin{choices}
  \item 10
  \item 50
  \item 65
  \item 80
  \item 100
\end{choices}
\end{question}
\answer{5}
